% from ega4-1-fr.tex
\oldpage[IV-1]{225}

suite, si $f:X\to Y$ est un $S$-morphisme de $S$-préschémas, et s'il existe
un recouvrement ouvert $(S_\lambda)$ de $S$ tel que les restrictions
$g^{-1}(S_\lambda)\to h^{-1}(S_\lambda)$ de $f$ (où $g$ et $h$ sont les
morphismes structuraux) soient quasi-compactes, alors $f$ est quasi-compact.

\begin{proposition}[1.1.2]
\label{IV.1.1.2-fr}
\begin{enumerate}
  \item[(i)] Une immersion $X\to Y$ est quasi-compacte si elle est fermée, ou si
  l'espace sous-jacent à $Y$ est localement noethérien ou si l'espace
  sous-jacent à $X$ est noethérien.
  \item[(ii)] Le composé de deux morphismes quasi-compacts est quasi-compact.
  \item[(iii)] Si $f:X\to Y$ est un $S$-morphisme quasi-compact, il en est de même de
  $f_{(S')}:X_{(S')}\to Y_{(S')}$ pour toute extension $g:S'\to S$ du
  préschéma de base.
  \item[(iv)] Si $f:X\to X'$ et $g:Y\to Y'$ sont deux $S$-morphismes
  quasi-compacts,
  \[
    f\times_S g:X\times_S Y\to X'\times_S Y'
  \]
  est quasi-compact.
  \item[(v)] Si le composé $g\circ f$ de deux morphismes $f:X\to Y$, $g:Y\to Z$
  est quasi-compact, et si $g$ est séparé, ou l'espace sous-jacent à $X$
  localement noethérien, $f$ est quasi-compact.
  \item[(vi)] Pour qu'un morphisme $f$ soit quasi-compact, il faut et il suffit que
  $f_{\mathrm{red}}$ le soit.
\end{enumerate}
\end{proposition}

Pour la démonstration, voir (I, 6.6.4). On notera que l'assertion (vi) est
aussi conséquence de la proposition plus générale~:

\begin{proposition}[1.1.3]
\label{IV.1.1.3-fr}
Soient $f:X\to Y$, $g:Y\to Z$ deux morphismes. Si $g\circ f$ est
quasi-compact et $f$ surjectif, $g$ est quasi-compact.
\end{proposition}

\begin{proof}
En effet, si $V$ est un ouvert quasi-compact dans $Z$,
$f^{-1}(g^{-1}(V))$ est quasi-compact par hypothèse, et l'on a
$g^{-1}(V)=f(f^{-1}(g^{-1}(V)))$, puisque $f$ est surjectif~; donc
$g^{-1}(V)$ est quasi-compact.
\end{proof}

\begin{corollary}[1.1.4]
\label{IV.1.1.4-fr}
Soient $f:X\to Y$, $g:Y'\to Y$ deux morphismes, $X'=X\times_Y Y'$,
$f'=f_{(Y')}:X'\to Y'$~; supposons $g$ quasi-compact et surjectif (resp.
surjectif). Alors, pour que $f$ soit quasi-compact (resp. dominant), il suffit
que $f'$ le soit.
\end{corollary}

\begin{proof}
Si $g':X'\to X$ est la projection canonique, $g'$ est un morphisme surjectif
(I, 3.5.2), et l'on a $f\circ g'=g\circ f'$~; si $f'$ est quasi-compact
(resp. dominant), il en est de même de $g\circ f'$ puisque $g$ est
quasi-compact (1.1.2) (resp. surjectif)~; comme $f\circ g'$ est quasi-compact
(resp. dominant) et que $g'$ est surjectif, on en déduit que $f$ est
quasi-compact (1.1.3) (resp. dominant).
\end{proof}

